\Article{6}
Bezeichnet \(p'+iq'\) eine beliebige fest gewählte Funktion von \(p+iq\) (wobei \(p',q'\) reelle Funktionen von \(p,q\) sind), so sieht man leicht, dass auch
\[
p'+iq'=\text{Const.}\quad\text{und}\quad p'-iq'=\text{Const.}
\]
\noindent die Integrale der Differentialgleichung \(\omega=0\) darstellen. Tatsächlich sind diese jeweils mit den obigen Gleichungen
\[
p+iq=\text{Const.}\quad\text{und}\quad p-iq=\text{Const.}
\]
\noindent völlig gleichbedeutend. Ebenso sind die Integrale der Differentialgleichung \(\Omega=0\)
\[
P'+iQ'=\text{Const.}\quad\text{und}\quad P'-iQ'=\text{Const.}
\]
\noindent mit den obigen Gleichungen
\[
P+iQ=\text{Const.}\quad\text{und}\quad P-iQ=\text{Const.}
\]
\noindent völlig gleichbedeutend, wenn \(P'+iQ'\) eine beliebige fest gewählte Funktion von \(P+iQ\) bezeichnet (wobei \(P',Q'\) reelle Funktionen von \(P,Q\) sind). Daraus folgt, dass in der allgemeinen Lösung unserer Aufgabe, die wir im vorhergehenden Artikel gegeben haben, auch \(p',q'\) an die Stelle von \(p,q\) und entsprechend \(P',Q'\) an die Stelle von \(P,Q\) treten können. Eine solche Änderung macht die Lösung zwar nicht allgemeiner, doch kann für die Anwendung bisweilen die eine Form für den einen und die andere für einen anderen Zweck bequemer sein.
\SourceLeaf{204}
\Article{7}
Bezeichnet man die Funktionen, die durch Differentiation der beliebigen Funktionen \(f,f'\) entstehen, mit \(\varphi\) beziehungsweise \(\varphi'\), sodass \(\mathrm{d}\mathord{.}f\upsilon=\varphi\upsilon\mathord{.}\mathrm{d}\upsilon\), \(\mathrm{d}\mathord{.}f'\upsilon=\varphi'\upsilon\mathord{.}\mathrm{d}\upsilon\), dann folgt aus unserer allgemeinen Lösung
\[
\frac{\mathrm{d}P+i\,\mathrm{d}Q}{\mathrm{d}p+i\,\mathrm{d}q}=\varphi(p+iq),\qquad
\frac{\mathrm{d}P-i\,\mathrm{d}Q}{\mathrm{d}p-i\,\mathrm{d}q}=\varphi'(p-iq)
\]
\noindent also
\[
\frac{mmn}{N}=\varphi(p+iq)\mathord{.}\varphi'(p-iq)
\]
\noindent Das Vergrößerungsverhältnis bestimmt sich daher durch die Formel
\[
m=\BareRoot{\Bigl\{\frac{\mathrm{d}p^2+\mathrm{d}q^2}{\omega}\mathord{.}\frac{\Omega}{\mathrm{d}P^2+\mathrm{d}Q^2}\mathord{.}\varphi(p+iq)\mathord{.}\varphi'(p-iq)\Bigr\}}
\]

